\documentclass{amsart}
\usepackage{amsmath,amssymb,amsthm,hyperref}
\usepackage{algorithm}
\usepackage{algpseudocode}

\newtheorem{theorem}{Theorem}
\newtheorem{lemma}{Lemma}
\newtheorem{proposition}{Proposition}
\newtheorem{corollary}{Corollary}
\theoremstyle{definition}
\newtheorem{definition}{Definition}
\theoremstyle{remark}
\newtheorem{remark}{Remark}

\newcommand{\rpo}{>_{\mathrm{rpo}}}
\newcommand{\rpoe}{\ge_{\mathrm{rpo}}}
\newcommand{\lpo}{>_{\mathrm{lpo}}}
\newcommand{\lpoe}{\ge_{\mathrm{lpo}}}
\newcommand{\sig}{\Sigma}
\newcommand{\prc}{>_{\sig}}
\newcommand{\T}{T(\sig,V)}
\newcommand{\size}[1]{|#1|}

\begin{document}

\title{Complexity of RPO and LPO comparison with a given (partial) precedence}
\maketitle

\section{Statement and main result}

Let $\sig$ be a finite signature whose function symbols have fixed arities, let
$V$ be a countable set of variables, and let $\T$ be the set of first-order terms
over $\sig$ and $V$. A \emph{precedence} $\prc$ is a strict partial order on
$\sig$. For a term $u$ we write $\size{u}$ for its size, i.e.\ the number of
nodes (occurrences of function symbols and variables) in its tree. The two
relations $\rpo$ and $\lpo$ are defined as in the problem statement. We make the
standard convention, used implicitly in the definition through the symbol
$\rpoe$ (resp.\ $\lpoe$), that $\rpoe$ is the reflexive closure of $\rpo$
\emph{together with} the equivalence generated by permuting multiset arguments;
we make this completely precise in Section~\ref{sec:equiv} and prove that under
either standard reading the answer is the same.

We consider the two decision problems:
\begin{description}
\item[(i) {\sc RPO-Compare}] Given $\sig$, a precedence $\prc$ (a finite list of
the pairs $f \prc g$), and $s,t\in\T$, decide whether $s \rpo t$.
\item[(ii) {\sc LPO-Compare}] Given $\sig$, a precedence $\prc$, and $s,t\in\T$,
decide whether $s \lpo t$.
\end{description}

\begin{theorem}\label{thm:main}
Both {\sc RPO-Compare} and {\sc LPO-Compare} are decidable in polynomial time.
Concretely, $s\rpo t$ and $s\lpo t$ can each be decided in time
$O\bigl((\size{s}\cdot\size{t})\cdot p(\size{s}+\size{t}+\size{\sig})\bigr)$ for
a fixed polynomial $p$; in particular both problems lie in $\mathrm{P}$ (and
hence in $\mathrm{NP}\cap\mathrm{coNP}\subseteq\mathrm{PSPACE}$). They are not
$\mathrm{NP}$-hard, $\mathrm{coNP}$-hard, or $\mathrm{PSPACE}$-hard unless
$\mathrm{P}$ coincides with the corresponding class.
\end{theorem}

The polynomiality holds even though the precedence is allowed to be a
\emph{non-total} partial order, and even though terms may contain variables.
This contrasts with the well-known fact (Krishnamoorthy and Narendran, 1985)
that the \emph{orientation} problem --- ``does there \emph{exist} a precedence
$\prc$ under which all rules of a given rewrite system decrease in $\rpo$ (resp.\
$\lpo$)?'' --- is $\mathrm{NP}$-complete; see Remark~\ref{rem:existence}. In the
present problems the precedence is part of the input, which is what makes the
comparison tractable.

The whole difficulty of the multiset case lies in the multiset extension of a
\emph{partial} order modulo a nontrivial equivalence. The key technical fact
(Lemma~\ref{lem:mul}) is that, because the equivalence arising in RPO is a
\emph{congruence} for $\rpo$, this multiset comparison reduces to one maximum
bipartite matching followed by a ``for-all/exists'' check, and hence is
polynomial; without the congruence property the multiset comparison of a partial
order modulo an equivalence is genuinely intractable, and it is precisely the
congruence property that rescues tractability here.

\section{Preliminaries on terms and on the orders}\label{sec:equiv}

Write a term either as a variable $x\in V$, or as $f(u_1,\dots,u_k)$ with
$f\in\sig$ of arity $k$. The \emph{subterms} of $u$, denoted $\mathrm{Sub}(u)$,
are $u$ together with the subterms of its immediate arguments. The number of
distinct subterm \emph{values} of $u$ is at most $\size{u}$. A variable $x$
\emph{occurs strictly} in $u$ if $x$ occurs in some immediate argument of $u$ (in
particular $u$ is not itself the variable).

\subsection{The equivalence and the quasi-order}
For RPO with multiset status, define the relation $\sim_{\mathrm{rpo}}$ on $\T$
recursively as the least relation with: $x\sim_{\mathrm{rpo}}y$ iff $x=y$ for
variables; and $f(s_1,\dots,s_m)\sim_{\mathrm{rpo}}g(t_1,\dots,t_n)$ iff $f=g$
(so $m=n$) and there is a permutation $\pi$ of $\{1,\dots,m\}$ with
$s_i\sim_{\mathrm{rpo}}t_{\pi(i)}$ for all $i$. We set
$s\rpoe t$ iff $s\rpo t$ or $s\sim_{\mathrm{rpo}}t$.

For LPO (lexicographic status) we use $s\sim_{\mathrm{lpo}}t$ iff $s=t$
(syntactic equality), and $s\lpoe t$ iff $s\lpo t$ or $s=t$; here no nontrivial
permutations occur because the lexicographic extension compares argument tuples
positionally.

\begin{remark}[Robustness of the reading of $\rpoe$]\label{rem:reading}
Two readings of $\rpoe$ are common: (a) the reflexive closure
$s\rpo t \text{ or } s=t$ (syntactic equality), and (b) the quasi-order
$s\rpo t \text{ or } s\sim_{\mathrm{rpo}}t$ above. Reading (b) is the standard
one for the multiset path order, and is the one we use; reading (a) is a special
case obtained by forbidding the permutation step in $\sim_{\mathrm{rpo}}$. The
algorithm and proof below handle reading (b) (the more general and standard
one); the polynomial bound holds verbatim for reading (a) as well, since (a)
only removes work. Thus Theorem~\ref{thm:main} holds under either reading.
\end{remark}

\subsection{The multiset extension as a congruence}
The next two lemmas isolate the two structural facts we need about
$\sim_{\mathrm{rpo}}$ and the multiset comparison.

\begin{lemma}[Congruence]\label{lem:cong}
$\sim_{\mathrm{rpo}}$ is an equivalence relation, and it is a congruence for
$\rpo$: if $a\sim_{\mathrm{rpo}}a'$ and $b\sim_{\mathrm{rpo}}b'$, then
$a\rpo b\iff a'\rpo b'$. In particular, for every term $u$, the set of terms
$\rpo$-below $u$ and the set of terms $\rpo$-above $u$ are unions of
$\sim_{\mathrm{rpo}}$-classes.
\end{lemma}

\begin{proof}
Reflexivity, symmetry and transitivity of $\sim_{\mathrm{rpo}}$ follow by
structural induction (identity permutation; inverse permutation; composition of
permutations). For the congruence property we argue by induction on
$\size{a}+\size{b}$ that $a\sim_{\mathrm{rpo}}a'$ and $b\sim_{\mathrm{rpo}}b'$
imply $a\rpo b\Rightarrow a'\rpo b'$ (the converse is symmetric, and the
equivalence then follows). If $a$ is a variable then $a\rpo b$ never holds, so
the implication is vacuous; thus $a=f(a_1,\dots,a_m)$ and, since
$a\sim_{\mathrm{rpo}}a'$, $a'=f(a'_1,\dots,a'_m)$ with a permutation $\rho$ and
$a_i\sim_{\mathrm{rpo}}a'_{\rho(i)}$.

Suppose $a\rpo b$.
\emph{Case $b$ a variable.} Then $b'=b$ (variables are $\sim_{\mathrm{rpo}}$ only
to themselves) and $a\rpo b$ holds by clause~(1) because $b$ occurs strictly in
$a$. The variables occurring strictly in $a$ and in $a'$ coincide as sets,
because $a$ and $a'$ have the same multiset of immediate arguments up to
$\sim_{\mathrm{rpo}}$, and $\sim_{\mathrm{rpo}}$-equivalent terms contain exactly
the same variables (immediate from the definition of $\sim_{\mathrm{rpo}}$ by
induction). Hence $b$ occurs strictly in $a'$ and $a'\rpo b'$ by~(1).

\emph{Case $b=g(b_1,\dots,b_n)$.} Then $b'=g(b'_1,\dots,b'_n)$ with a permutation
$\sigma$, $b_j\sim_{\mathrm{rpo}}b'_{\sigma(j)}$. We check that whichever of
clauses (2a)--(2c) certifies $a\rpo b$ also certifies $a'\rpo b'$.
\begin{itemize}
\item (2a): $a_i\rpoe b$ for some $i$. Then $a'_{\rho(i)}\sim_{\mathrm{rpo}}a_i$,
and by the induction hypothesis (with the smaller term $a_i$) together with
$b\sim_{\mathrm{rpo}}b'$ we get $a'_{\rho(i)}\rpoe b'$ (if $a_i\rpo b$ use the
IH; if $a_i\sim_{\mathrm{rpo}}b$ then $a'_{\rho(i)}\sim_{\mathrm{rpo}}b'$ by
transitivity). Hence (2a) holds for $a',b'$.
\item (2b): $f\prc g$ and $a\rpo b_j$ for all $j$. By the IH (smaller second
argument $b_j$, using $a\sim_{\mathrm{rpo}}a'$ and
$b_j\sim_{\mathrm{rpo}}b'_{\sigma(j)}$) we get $a'\rpo b'_{\sigma(j)}$ for all
$j$, i.e.\ $a'\rpo b'_{j'}$ for all $j'$, and $f\prc g$ is unchanged. Hence (2b)
holds.
\item (2c): $f=g$ and $\{a_1,\dots,a_m\}\rpo^{\mathrm{mul}}\{b_1,\dots,b_n\}$.
Since the multisets $\{a'_i\}$ and $\{a_i\}$ agree up to $\sim_{\mathrm{rpo}}$,
and likewise $\{b'_j\}$ and $\{b_j\}$, and since (by the invariance statement of
Lemma~\ref{lem:mul} below, whose proof only uses $\rpo$ and $\sim_{\mathrm{rpo}}$
on strictly smaller terms) the multiset order is invariant under replacing
elements by $\sim_{\mathrm{rpo}}$-equivalent ones, we get
$\{a'_i\}\rpo^{\mathrm{mul}}\{b'_j\}$. Hence (2c) holds.
\end{itemize}
In all cases $a'\rpo b'$, completing the induction. (No circularity arises:
clause~(2c) and Lemma~\ref{lem:mul} are applied only to the immediate arguments,
which have strictly smaller size, so the inductive hypothesis governs them.)
\end{proof}

\section{The multiset comparison is polynomial}

We make precise the multiset extension and show how to compute it. Let $(P,>)$ be
a strict partial order equipped with an equivalence $\equiv$ that is a congruence
for $>$ (i.e.\ $x\equiv x'$, $y\equiv y'\Rightarrow(x>y\iff x'>y')$ and
$x\equiv y\Rightarrow x\not> y$). For finite multisets $M,N$ over $P$, the
\emph{multiset extension} $M\,{>}^{\mathrm{mul}}\,N$ is, by the
Dershowitz--Manna definition,
\[
M\,{>}^{\mathrm{mul}}\,N \iff
\exists\, X,Y \text{ finite multisets over } P:\;
X\neq\emptyset,\;\;
N \equiv (M\setminus X)\uplus Y,\;\;
\forall y\in Y\;\exists x\in X:\; x>y,
\tag{$\star$}
\]
where $\setminus$ and $\uplus$ are multiset difference and union, and the
relation $N\equiv(M\setminus X)\uplus Y$ between multisets means that there is an
occurrence-wise bijection matching $\equiv$-equivalent elements.

\begin{lemma}[Polynomial multiset comparison]\label{lem:mul}
Let $(P,>,\equiv)$ be as above with $>$ and $\equiv$ decidable, and let $M,N$ be
multisets with $|M|+|N|=k$. Then:
\begin{enumerate}
\item[\textup{(a)}] (Invariance) $M\,{>}^{\mathrm{mul}}\,N$ is unchanged if any
element of $M$ or $N$ is replaced by an $\equiv$-equivalent one.
\item[\textup{(b)}] (Algorithm) $M\,{>}^{\mathrm{mul}}\,N$ holds iff the following
procedure outputs ``yes''. \emph{Step 1:} compute a maximum bipartite matching
$R$ between the occurrences of $M$ and of $N$ with edge relation $\equiv$.
\emph{Step 2:} let $M'$ (resp.\ $N'$) be the $R$-unmatched occurrences of $M$
(resp.\ $N$). \emph{Step 3:} output ``yes'' iff $M'\neq\emptyset$ and for every
$y\in N'$ there is $x\in M'$ with $x>y$. The procedure runs in time
$O(k^{3}+k^{2}c)$, where $c$ bounds one evaluation of $>$ or $\equiv$.
\end{enumerate}
\end{lemma}

\begin{proof}
\emph{(a)} The predicate $(\star)$ refers only to $>$ and to $\equiv$-matchings.
Replacing an element by an $\equiv$-equivalent one preserves all $\equiv$-edges
and, by the congruence property, all relations $x>y$; hence $(\star)$ is
unchanged.

\emph{(b)} By part (a) we may pass to the quotient $P/\!\equiv$: define
$\bar P=P/\!\equiv$, with $[x]>[y]$ iff $x>y$ (well defined and a strict partial
order because $\equiv$ is a congruence for $>$ and $x\equiv y\Rightarrow x\not>y$),
and let $\bar M,\bar N$ be the images of $M,N$. Then
$M\,{>}^{\mathrm{mul}}\,N$ over $(P,>,\equiv)$ holds iff
$\bar M\,{>}^{\mathrm{mul}}\,\bar N$ over $(\bar P,>)$ with $\equiv$ now equality:
indeed $(\star)$ only depends on the data preserved by quotienting. Over a strict
partial order with equality, the Dershowitz--Manna order $(\star)$ has the
classical equivalent characterization \cite{DM}
\[
\bar M\,{>}^{\mathrm{mul}}\,\bar N \iff
\bar M\neq\bar N \ \text{and}\
\forall\, y\in \bar N\setminus\bar M\ \exists\, x\in \bar M\setminus\bar N:\ x>y,
\tag{$\star\star$}
\]
where the differences are multiset differences with respect to equality.

It remains to check that the procedure in (b), run on the original $M,N$, decides
$(\star\star)$ for $\bar M,\bar N$. A maximum $\equiv$-matching $R$ of $M$ and $N$
projects to a maximum equality-matching of $\bar M$ and $\bar N$; the maximum
equality-matching of two multisets removes exactly their common part, so the
projected leftovers are $\bar M\setminus\bar N$ and $\bar N\setminus\bar M$. Thus
$\overline{M'}=\bar M\setminus\bar N$ and $\overline{N'}=\bar N\setminus\bar M$.
Now:
\begin{itemize}
\item If the procedure outputs ``yes'', then $M'\neq\emptyset$ gives
$\bar M\neq\bar N$ (a maximum matching leaves $M$ entirely matched iff
$\bar M\subseteq\bar N$ as multisets, which together with $|M|=|M'|+|R|$ forces
$\bar M=\bar N$ only when also $N'=\emptyset$; if $M'\neq\emptyset$ then
$\bar M\setminus\bar N\neq\emptyset$, so $\bar M\neq\bar N$), and the domination
test for $M',N'$ is exactly the right-hand condition of $(\star\star)$ for
$\bar M\setminus\bar N,\ \bar N\setminus\bar M$. Hence $(\star\star)$ holds, so
$M\,{>}^{\mathrm{mul}}\,N$.
\item Conversely, if $(\star\star)$ holds for $\bar M,\bar N$, then
$\bar M\setminus\bar N\neq\emptyset$ (otherwise $\bar M$ is a sub-multiset of
$\bar N$; if also $\bar N\setminus\bar M=\emptyset$ then $\bar M=\bar N$,
contradicting $\bar M\neq\bar N$; if $\bar N\setminus\bar M\neq\emptyset$ then its
elements cannot be dominated from the empty set $\bar M\setminus\bar N$,
contradicting $(\star\star)$). Therefore $\overline{M'}=\bar M\setminus\bar N\neq
\emptyset$, so $M'\neq\emptyset$; and the domination condition holds by
$(\star\star)$. Hence the procedure outputs ``yes''.
\end{itemize}
Thus the procedure correctly decides $M\,{>}^{\mathrm{mul}}\,N$.

\emph{Running time.} Building the $\equiv$-edge set costs $O(k^{2})$ evaluations
of $\equiv$; maximum bipartite matching by augmenting paths on $\le k$ vertices
per side and $\le k^{2}$ edges costs $O(k^{3})$; the domination test costs
$O(k^{2})$ evaluations of $>$. Total $O(k^{3}+k^{2}c)$.
\end{proof}

\begin{remark}[Why partiality alone does not break tractability]
The equivalence $(\star)\Leftrightarrow(\star\star)$ holds for \emph{every}
strict partial order, not only total ones; this is the content of \cite{DM}. The
only potential source of intractability would be a nontrivial equivalence
$\equiv$ that is \emph{not} a congruence for $>$: in that case the choice of
\emph{which} maximum matching to remove can matter and the comparison becomes a
hard assignment problem. Lemma~\ref{lem:cong} guarantees that for RPO the
relevant equivalence $\sim_{\mathrm{rpo}}$ \emph{is} a congruence for $\rpo$, so
this obstruction never arises and Lemma~\ref{lem:mul} applies. This is the exact
point where ``RPO comparison with a given partial precedence'' stays in
$\mathrm{P}$.
\end{remark}

\section{Polynomial-time algorithms}

\subsection{RPO}
Fix the input $\sig,\prc,s,t$. We tabulate, for $a\in\mathrm{Sub}(s)$ and
$b\in\mathrm{Sub}(t)$, the predicates
$\textsc{Eq}(a,b)\equiv[a\sim_{\mathrm{rpo}}b]$ and
$\textsc{Gt}(a,b)\equiv[a\rpo b]$ (and $\textsc{Ge}=\textsc{Gt}\lor\textsc{Eq}$).
The crucial closure property is:

\begin{lemma}[Subterm closure of the recursion]\label{lem:closure}
In evaluating $s\rpo t$ by the defining clauses, every comparison $a\rpo b$ that
is ever invoked has $a\in\mathrm{Sub}(s)$ and $b\in\mathrm{Sub}(t)$; likewise
every $a\sim_{\mathrm{rpo}}b$ invoked has $a\in\mathrm{Sub}(s)$,
$b\in\mathrm{Sub}(t)$. Moreover each such comparison depends only on comparisons
of strictly smaller total size $\size{a}+\size{b}$.
\end{lemma}

\begin{proof}
Consider $a=f(a_1,\dots,a_p)\rpo b=g(b_1,\dots,b_q)$ with $a\in\mathrm{Sub}(s)$,
$b\in\mathrm{Sub}(t)$:
\begin{itemize}
\item (2a) invokes $a_i\rpoe b$: $a_i\in\mathrm{Sub}(a)\subseteq\mathrm{Sub}(s)$,
$b\in\mathrm{Sub}(t)$, strictly smaller first size, hence strictly smaller total
size.
\item (2b) invokes $a\rpo b_j$: $a\in\mathrm{Sub}(s)$,
$b_j\in\mathrm{Sub}(b)\subseteq\mathrm{Sub}(t)$, strictly smaller second size.
\item (2c) invokes (via Lemma~\ref{lem:mul}) only $a_i\rpo b_j$ and
$a_i\sim_{\mathrm{rpo}}b_j$ on immediate arguments: both in $\mathrm{Sub}(s)$,
$\mathrm{Sub}(t)$ respectively, strictly smaller total size.
\end{itemize}
If $a$ is a variable, $a\rpo b$ is false and there is no recursive call. The
definition of $\sim_{\mathrm{rpo}}$ likewise calls $a_i\sim_{\mathrm{rpo}}b_j$
only on immediate arguments. Since at the top $a=s\in\mathrm{Sub}(s)$,
$b=t\in\mathrm{Sub}(t)$, and each bullet keeps the first argument in
$\mathrm{Sub}(s)$ and the second in $\mathrm{Sub}(t)$, the property is preserved
at every level.
\end{proof}

\begin{algorithm}
\caption{{\sc RPO-Compare}: decide $s\rpo t$}
\begin{algorithmic}[1]
\Require Signature $\sig$, precedence $\prc$ (as a queryable relation), terms $s,t$
\Ensure \textbf{true} iff $s\rpo t$
\State Compute the distinct subterm sets $A=\mathrm{Sub}(s)$, $B=\mathrm{Sub}(t)$
\State Sort all pairs $(a,b)\in A\times B$ by nondecreasing $\size{a}+\size{b}$
\For{each pair $(a,b)$ in this order}
  \State \Comment{$\textsc{Eq}(a,b)=[a\sim_{\mathrm{rpo}}b]$}
  \If{$a,b$ are the same variable} $\textsc{Eq}(a,b)\gets$ \textbf{true}
  \ElsIf{$a,b$ have the same top symbol $f$ of arity $k$}
     \State $\textsc{Eq}(a,b)\gets$ [there is a perfect matching of the immediate
     arguments of $a$ and $b$ under the already-computed relation $\textsc{Eq}$]
  \Else\ $\textsc{Eq}(a,b)\gets$ \textbf{false}
  \EndIf
  \State \Comment{$\textsc{Gt}(a,b)=[a\rpo b]$}
  \If{$a$ is a variable} $\textsc{Gt}(a,b)\gets$ \textbf{false}
  \ElsIf{$b$ is a variable} $\textsc{Gt}(a,b)\gets$ [$b$ occurs strictly in $a$]
  \Else
     \State write $a=f(a_1,\dots,a_p)$, $b=g(b_1,\dots,b_q)$
     \State $\mathit{c2a}\gets \exists i:\ \textsc{Gt}(a_i,b)\lor\textsc{Eq}(a_i,b)$
     \State $\mathit{c2b}\gets (f\prc g)\ \wedge\ \forall j:\ \textsc{Gt}(a,b_j)$
     \State $\mathit{c2c}\gets (f=g)\ \wedge\ \textsc{MulGt}\bigl(\{a_1,\dots,a_p\},\{b_1,\dots,b_q\}\bigr)$
        \Comment{Lemma~\ref{lem:mul}, using $\textsc{Gt},\textsc{Eq}$ on immediate args}
     \State $\textsc{Gt}(a,b)\gets \mathit{c2a}\lor \mathit{c2b}\lor \mathit{c2c}$
  \EndIf
\EndFor
\State \Return $\textsc{Gt}(s,t)$
\end{algorithmic}
\end{algorithm}

\begin{proposition}[Correctness and complexity of {\sc RPO-Compare}]\label{prop:rpo}
The algorithm outputs \textbf{true} iff $s\rpo t$, and runs in time
$O\bigl(\size{s}^2\,\size{t}^2+\size{s}\,\size{t}\,\size{\sig}^2\bigr)$. Hence
{\sc RPO-Compare}$\in\mathrm{P}$.
\end{proposition}

\begin{proof}
\emph{Well-foundedness of the schedule.} By Lemma~\ref{lem:closure},
$\textsc{Gt}(a,b)$ depends only on values $\textsc{Gt}(a',b'),\textsc{Eq}(a',b')$
with $\size{a'}+\size{b'}<\size{a}+\size{b}$, and $\textsc{Eq}(a,b)$ depends only
on $\textsc{Eq}$ of immediate arguments, which have strictly smaller size.
Processing pairs in nondecreasing $\size{a}+\size{b}$ ensures every referenced
value is already computed.

\emph{Correctness.} By induction on $\size{a}+\size{b}$, each assignment
implements its defining clause exactly. The $\textsc{Eq}$ branch implements
$\sim_{\mathrm{rpo}}$: a perfect matching of the immediate arguments under
$\textsc{Eq}$ is precisely the required permutation $\pi$. The $\textsc{Gt}$
branch implements clauses (1), (2a), (2b), (2c) verbatim, with (2c) computed
correctly by Lemma~\ref{lem:mul}(b), which is applicable because
$\sim_{\mathrm{rpo}}$ is a congruence for $\rpo$ (Lemma~\ref{lem:cong}). Finally,
$\rpo$ is by definition the least relation closed under these monotone clauses,
and Lemma~\ref{lem:closure} shows the dependency relation among the comparisons
is well-founded (strictly decreasing total size); a least relation closed under
monotone clauses with a well-founded dependency relation is computed exactly by
bottom-up evaluation along that order. Hence $\textsc{Gt}(s,t)=[s\rpo t]$.

\emph{Complexity.} There are at most $\size{s}\cdot\size{t}$ pairs. For each,
$\textsc{Eq}$ costs one bipartite matching on $\le K$ vertices per side ($K$ the
maximum arity), i.e.\ $O(K^{3})$ table look-ups; $\textsc{Gt}$ costs $O(p+q)$
look-ups for (2a),(2b) plus, for (2c), one $\textsc{MulGt}$ costing $O((p+q)^{3})$
look-ups by Lemma~\ref{lem:mul}(b). Each precedence query $f\prc g$ costs
$O(1)$ after an $O(\size{\sig}^2)$ preprocessing of the partial order. Summing
gives $O\bigl(\size{s}\,\size{t}\,(K^{3})+\size{\sig}^2\bigr)$; since
$K\le\min(\size{s},\size{t})$ this is bounded by
$O\bigl(\size{s}^2\size{t}^2+\size{s}\,\size{t}\,\size{\sig}^2\bigr)$, polynomial
in the input size. Thus {\sc RPO-Compare}$\in\mathrm{P}$.
\end{proof}

\subsection{LPO}
The LPO algorithm is identical except that clause (2c) is the lexicographic test
and the equivalence is syntactic equality.

\begin{algorithm}
\caption{{\sc LPO-Compare}: decide $s\lpo t$}
\begin{algorithmic}[1]
\Require Signature $\sig$, precedence $\prc$, terms $s,t$
\Ensure \textbf{true} iff $s\lpo t$
\State Compute $A=\mathrm{Sub}(s)$, $B=\mathrm{Sub}(t)$
\For{each $(a,b)\in A\times B$ in nondecreasing order of $\size{a}+\size{b}$}
  \If{$a$ is a variable} $\textsc{Gt}(a,b)\gets$ \textbf{false}
  \ElsIf{$b$ is a variable} $\textsc{Gt}(a,b)\gets$ [$b$ occurs strictly in $a$]
  \Else
     \State $a=f(a_1,\dots,a_p)$, $b=g(b_1,\dots,b_q)$
     \State $\mathit{c2a}\gets \exists i:\ \textsc{Gt}(a_i,b)\lor (a_i=b)$
     \State $\mathit{c2b}\gets (f\prc g)\ \wedge\ \forall j:\ \textsc{Gt}(a,b_j)$
     \State $\mathit{c2c}\gets (f=g)\ \wedge\ \bigl(\forall j:\ \textsc{Gt}(a,b_j)\bigr)\ \wedge\ \textsc{LexGt}\bigl((a_1,\dots,a_p),(b_1,\dots,b_q)\bigr)$
     \State $\textsc{Gt}(a,b)\gets \mathit{c2a}\lor \mathit{c2b}\lor \mathit{c2c}$
  \EndIf
\EndFor
\State \Return $\textsc{Gt}(s,t)$
\end{algorithmic}
\end{algorithm}

Here $\textsc{LexGt}\bigl((a_1,\dots,a_p),(b_1,\dots,b_q)\bigr)$ scans
$i=1,2,\dots$: while $a_i=b_i$ continue; at the first index $i$ with $a_i\neq b_i$
return $\textsc{Gt}(a_i,b_i)$; if one tuple is a proper prefix of the other,
return \textbf{true} iff the first tuple is longer. (For $f=g$ the arities are
equal, $p=q$, so the prefix case does not occur; the general form is stated for
completeness.)

\begin{proposition}[Correctness and complexity of {\sc LPO-Compare}]\label{prop:lpo}
The algorithm outputs \textbf{true} iff $s\lpo t$, and runs in time
$O\bigl(\size{s}^2\size{t}^2+\size{s}\,\size{t}\,\size{\sig}^2\bigr)$. Hence
{\sc LPO-Compare}$\in\mathrm{P}$.
\end{proposition}

\begin{proof}
The subterm-closure lemma (Lemma~\ref{lem:closure}) holds verbatim: (2c) now
invokes $\textsc{Gt}(a,b_j)$ (second argument strictly smaller) and, via
$\textsc{LexGt}$, $\textsc{Gt}(a_i,b_i)$ with both arguments strictly smaller (so
strictly smaller total size). The lexicographic comparison uses syntactic
equality $a_i=b_i$, decided in $O(\size{a_i})$ by a tree traversal (or $O(1)$
after hash-consing the subterms, which costs $O(\size{s}+\size{t})$). Hence the
recursion is well-founded and bottom-up evaluation in nondecreasing
$\size{a}+\size{b}$ is correct, by the least-fixed-point argument of
Proposition~\ref{prop:rpo}. Each $\textsc{LexGt}$ scans at most
$\min(p,q)\le K$ positions with $O(1)$ table look-ups and one equality check per
position, so each pair costs $O(K+\size{\sig}^2)$ besides equality checks, and
the total is $O\bigl(\size{s}^2\size{t}^2+\size{s}\,\size{t}\,\size{\sig}^2\bigr)$.
Thus {\sc LPO-Compare}$\in\mathrm{P}$.
\end{proof}

\section{Conclusion of the classification}

\begin{proof}[Proof of Theorem~\ref{thm:main}]
Propositions~\ref{prop:rpo} and~\ref{prop:lpo} give deterministic
polynomial-time algorithms for {\sc RPO-Compare} and {\sc LPO-Compare}, with the
stated running times; hence both problems are in $\mathrm{P}$. Since
$\mathrm{P}\subseteq\mathrm{NP}\cap\mathrm{coNP}\subseteq\mathrm{PSPACE}$, both
lie in $\mathrm{NP}$, $\mathrm{coNP}$, and $\mathrm{PSPACE}$. If either problem
were $\mathrm{NP}$-hard (resp.\ $\mathrm{coNP}$-, $\mathrm{PSPACE}$-hard) under
polynomial-time reductions, then a polynomial-time algorithm for a hard problem
lying in $\mathrm{P}$ would collapse the class into $\mathrm{P}$, giving
$\mathrm{P}=\mathrm{NP}$ (resp.\ $\mathrm{P}=\mathrm{coNP}$,
$\mathrm{P}=\mathrm{PSPACE}$). Hence no such hardness holds unless the
corresponding collapse occurs, and the tight classification of both problems is:
\emph{decidable in deterministic polynomial time}.
\end{proof}

\begin{remark}[On lower bounds within $\mathrm{P}$]\label{rem:lower}
The membership in $\mathrm{P}$ is the tight classification with respect to the
hierarchy $\mathrm{P}\subseteq\mathrm{NP},\mathrm{coNP}\subseteq\mathrm{PSPACE}$
about which the problem asks. Within $\mathrm{P}$, the natural lower-bound
phenomenon is sequentiality: the defining clauses form an alternation of
conjunctions (``$\forall j$'') and disjunctions (over cases / ``$\exists i$''),
so the value $\textsc{Gt}(s,t)$ is the output of a monotone Boolean circuit whose
gates are the pairs $(a,b)$. This places both problems in the realm of
monotone-circuit evaluation. We do not assert $\mathrm{P}$-completeness, as the
requested classification is with respect to
$\mathrm{P}/\mathrm{NP}/\mathrm{coNP}/\mathrm{PSPACE}$, where the decisive answer
is: \emph{both problems are in $\mathrm{P}$}.
\end{remark}

\begin{remark}[Contrast with the existence/orientation problem]\label{rem:existence}
The $\mathrm{NP}$-completeness commonly associated with recursive path orderings
\cite{KN} concerns a \emph{different} problem: given a finite rewrite system
$\{l_i\to r_i\}$, decide whether there \emph{exists} a precedence $\prc$ making
$l_i\rpo r_i$ for all $i$. There the precedence is \emph{not} part of the input
and must be guessed, and the problem is $\mathrm{NP}$-complete (the analogous LPO
orientability problem is also $\mathrm{NP}$-complete). In the problems considered
here the precedence is given, removing the source of nondeterminism and yielding
the polynomial-time algorithms above. The reason the comparison stays polynomial
even with a non-total precedence is precisely Lemma~\ref{lem:cong}: the
equivalence $\sim_{\mathrm{rpo}}$ is a congruence for $\rpo$, so the multiset
comparison (Lemma~\ref{lem:mul}) reduces to one maximum matching plus a
domination test rather than to an intractable assignment problem.
\end{remark}

\begin{thebibliography}{9}
\bibitem{DM} N.~Dershowitz and Z.~Manna,
\emph{Proving termination with multiset orderings},
Communications of the ACM 22 (1979) 465--476.
\bibitem{KN} M.~S.~Krishnamoorthy and P.~Narendran,
\emph{On recursive path ordering},
Theoretical Computer Science 40 (1985) 323--328.
\end{thebibliography}

\end{document}
